\subsection{与一个右有界平坦复形作张量积}

引理1. --- 设C为右A-模复形，E为左A-模复形。假设H(C) = 0且E是平坦的且在右侧有界。则H(C ⊗\_A E) = 0。

对于 \(k \in \mathbf{Z}\) ，设 \(\mathbf{T}^{(k)}\) 为 \(\mathbf{C} \otimes_{\mathbf{A}} \mathbf{E}\) 的子复形，使得

\[
\mathrm{T}_{n}^{(k)} = \sum_{\substack{p + q = n\\ q\leqslant k}}\mathrm{C}_{p}\otimes_{\mathrm{A}}\mathrm{E}_{q};
\]

那么 \(\mathbf{T}^{(k - 1)}\subset \mathbf{T}^{(k)}\) 并且我们有一个复形的正合序列

\[
0 \longrightarrow \mathrm{T} ^ {(k - 1)} \stackrel {i _ {k}} {\longrightarrow} \mathrm{T} ^ {(k)} \stackrel {\pi} {\longrightarrow} \mathrm{C} \otimes_ {\mathbf {A}} \mathrm{E} _ {k} (- k) \longrightarrow 0
\]

其中 \(i_{k}\) 是典范单射，\(\pi\) 将前面的直和投影到其因子 \(\mathbf{C}_{n-k} \otimes_{\mathbf{A}} \mathbf{E}_{k} = (\mathbf{C} \otimes_{\mathbf{A}} \mathbf{E}_{k}(-k))_{n}\) 上。由上述推论2，我们有 \(\mathbf{H}(\mathbf{C} \otimes_{\mathbf{A}} \mathbf{E}_{k}(-k)) = 0\) ，因此 \(i_{k}\) 是一个拟同构。我们有 \(\mathbf{T}^{(k)} = 0\)（当 \(k\) 足够小时），因为E在右侧有界，因此对k归纳可得对所有k有 \(\mathrm{H}(\mathrm{T}^{(k)}) = 0\)。最后，典范态射 \(\lim_{\mathbf{A}} \mathrm{T}^{(k)} \to \mathbf{C} \otimes_{\mathbf{A}} \mathbf{E}\) 显然是一个同构，因此 \(\mathrm{H}(\mathbf{C} \otimes_{\mathbf{A}} \mathbf{E}) = 0\) (X，第28页，命题1)。

引理2. --- 若 \(u: \mathbf{C} \to \mathbf{C}'\) 是右A-模复形的一个态射，且E是左A-模复形，则复形Con \((u) \otimes_{\mathbf{A}} \mathbf{E}\) 以及 Con \((u \otimes 1_{\mathbf{E}})\) 是同构的。

根据定义，Con \((u) \otimes_{\mathbf{A}} \mathbf{E}\) 是分次模 \((\mathbf{C}'(-1) \oplus \mathbf{C}) \otimes_{\mathbf{A}} \mathbf{E}\) 赋予微分D，使得对于 \(x \in \mathbf{C}_p\) , \(y' \in \mathbf{C}'(-1)_p = \mathbf{C}_{p-1}'\) , \(z \in \mathbf{E}_q\) ，有

\[
\mathrm{D} ((y ^ {\prime}, x) \otimes z) = (- d y ^ {\prime}, d x - u (y ^ {\prime})) \otimes z + (- 1) ^ {p} (y ^ {\prime}, x) \otimes d z,\tag{4}
\]

而 \(\operatorname{Con}(u \otimes 1_{\mathbf{E}})\) 是分次模 \((\mathbf{C}' \otimes_{\mathbf{A}} \mathbf{E})(-1) \oplus (\mathbf{C} \otimes_{\mathbf{A}} \mathbf{E})\) 配备微分 \(\mathbf{D}_1\) 使得，对于 \(x \in \mathbf{C}_p\) , \(y' \in \mathbf{C}_{p-1}'\) , \(z \in \mathbf{E}_q\) ，有

\[
\mathrm{D} _ {1} \left(y ^ {\prime} \otimes z, x \otimes z\right) = (- d y ^ {\prime} \otimes z - (- 1) ^ {p - 1} y ^ {\prime} \otimes d z, d x \otimes z + (- 1) ^ {p} x \otimes d z - u \left(y ^ {\prime}\right) \otimes z)
\]

\[
= \left(- d y ^ {\prime} \otimes z, (d x - u (y ^ {\prime})) \otimes z\right) + (- 1) ^ {p} \left(y ^ {\prime} \otimes d z, x \otimes d z\right)
\]

由此得出该论断。

命题4. — 设 \(u: \mathbf{C} \to \mathbf{C}'\) 是右A-模复形间的一个拟同构，且E是左A-模的一个复形，平坦且在右侧有界。则

\[
u \otimes 1 _ {\mathbf {E}}: \mathbf {C} \otimes_ {\mathbf {A}} \mathbf {E} \rightarrow \mathbf {C} ^ {\prime} \otimes_ {\mathbf {A}} \mathbf {E}
\]

是k-模复形间的一个拟同构。

事实上，根据X，第38页，推论， \(u\) (相应地， \(u \otimes 1_{\mathrm{E}}\) ）是拟同构当且仅当 \(\mathrm{H}(\mathrm{Con}(u)) = 0\) (相应地， \(\mathrm{H}(\mathrm{Con}(u \otimes 1_{\mathrm{E}})) = 0\) ）。然后我们由引理1和引理2得出结论。

注记。--- 利用交换同构，由前面的陈述可推出将张量积的两个参数角色互换后所得的类似陈述。

